\section{Introduction}
\label{sec:introduction}

This report investigates how pretrained models can be adapted to new tasks and how additional computation can improve their predictions. First, we fine-tune a language model on mathematical solutions using full fine-tuning, starting-point regularization, LoRA, and OPLoRA, comparing mathematical improvement with retention of existing capabilities and resource use \citep{li2018explicit,hu2021lora,xiong2025oplora,biderman2024lora}. Next, we connect a pretrained vision encoder to a language model and train it through caption alignment and visual question answering, using complementary image pairs and image controls to examine whether its answers rely on relevant visual evidence \citep{liu2023visual,goyal2017making}. Finally, we study test-time computation through repeated sampling, confidence-based selection, answer consensus, longer responses, sequential refinement, and adaptive sampling \citep{wang2022selfconsistency,kang2025selfcertainty,usman2026visconf,madaan2023selfrefine,aggarwal2023adaptive}. Across all three tasks, controlled comparisons and failure cases help assess what each method achieves, where it falls short, and how its benefits relate to computational cost.
